\documentclass[10pt,a4paper]{amsart}
\usepackage[margin=19mm]{geometry}
\usepackage{amsmath,amssymb,mathtools}
\usepackage[scr=rsfs,cal=euler]{mathalfa}
\usepackage{xfrac}
\usepackage[unicode,hidelinks,bookmarks=false]{hyperref}
\newcommand{\ie}{i.e.\ }
\newcommand{\cf}{cf.\ }
\newcommand{\resp}{respectively\ }
\newcommand{\Subsection}{\subsection}
\providecommand{\nobreakdash}{\nobreak\hbox{-}}
\setlength{\emergencystretch}{2em}
\multlinegap0pt
\usepackage{bm}
\usepackage{bbm}
\usepackage{dsfont}
\usepackage{xspace}
\usepackage{cancel}
\usepackage{mathtools}
\usepackage{xcolor}
\usepackage{enumitem}
\usepackage{amsthm}
\makeatletter
\@ifundefined{thm}{\newtheorem{thm}{Thm}}{}
\@ifundefined{lemma}{\newtheorem{lemma}[thm]{Lemma}}{}
\makeatother
\newcommand{\A}{\mathbb{A}}
\newcommand{\R}{\mathbb{R}}
\newcommand{\T}{\mathbb{T}}
\title[Conditions implying annular chaos]{Conditions implying annular chaos}
\author{Alejandro Passeggi, Fabio Armando Tal}
\date{Source: arXiv:2305.02963v4 (CC BY 4.0)}
\begin{document}
\maketitle

\noindent\emph{Source and licence.} Conditions implying annular chaos. Version arXiv:2305.02963v4 (CC BY 4.0), distributed under the Creative Commons Attribution 4.0 International licence, as recorded in the arXiv licence metadata for this exact version and shown on the abstract page https://arxiv.org/abs/2305.02963v4. Full rights notice: licenses/arxiv-2305.02963-CC-BY-4.0.txt.\par\smallskip

\noindent\textit{Editorial scope.} Counted unit: the Lemma whose source label is \texttt{casocurva} in the article (source0.tex lines 918--921, source label \texttt{casocurva}), delivered together with the article's own complete original proof of it (source0.tex lines 922--974, one \texttt{proof} environment of 53 lines). No mathematics is written, repaired or re-derived: statement and proof are the article's own text line for line. Required context retained uncounted: none. Every definition, display and label of those lines is retained unchanged. Omitted from this package and not redistributed: 1--917, 975--1602. Nothing inside a retained excerpt is omitted or altered: every retained body is delivered exactly as the article's lines are written, so no omission receipt of any kind accompanies this package. Citations: the retained excerpts carry no citation command at all, so no bibliography entry of the article is transcribed and no reference is redistributed. Reference adaptation: none was needed and none was made; every reference inside the delivered unit resolves inside it, so no unresolved reference or citation is produced. Printed slips kept literal: no printed word, symbol, hypothesis or number of the counted statement or of its proof was altered. Layout and class: the article's own class is \texttt{amsart} with packages including amssymb, amscd, amsthm, mathrsfs, enumitem, color, transparent, xcolor, inputenc, epsfig, graphicx, psfrag, float; the wrapper is the lane's pinned \texttt{amsart} editorial wrapper at 10pt on A4, which restates the theorem-like environments the retained text needs and adds the article's own macro definitions that the retained text needs (\texttt{\textbackslash A}, \texttt{\textbackslash R}, \texttt{\textbackslash T}), with extra wrapper packages (bm, bbm, dsfont, xspace, cancel, mathtools, xcolor, enumitem). The delivered numbering is the unit's own. Assets: no retained span contains a figure environment, a graphics-inclusion command, a PDF-page inclusion, a table or a TikZ picture; no image file is copied into this package, the package declares no asset dependency and no image file is redistributed with it. Licence: version arXiv:2305.02963v4 of this article is distributed under the Creative Commons Attribution 4.0 International licence, as recorded in the arXiv licence metadata for this exact version and shown on the abstract page https://arxiv.org/abs/2305.02963v4. A full rights notice is delivered at licenses/arxiv-2305.02963-CC-BY-4.0.txt. Characters: every retained excerpt is pure ASCII, so the delivered engine, XeLaTeX with its default Latin Modern text font, reports no missing character.
\begin{lemma}\label{casocurva}
Let $f\in\mathrm{Homeo}_0(\A)$ and $x_0$, $x_1$ two fixed points of $f$ with different rotation numbers. Consider 
a simple closed curve $\gamma$ invariant by $f$. Assume $x_0$ and $x_1$ are in the same connected component of the complement of $\gamma$, and let $U_0,U_1$ form a $2$-dpn for $x_0,x_1$. Then $\gamma$ cannot intersect both $U_0$ and $U_1$.
\end{lemma}

\begin{proof}
Assume, for a contradiction, that $\gamma$ intersects both $U_0$ and $U_1$. We can assume, by a change of coordinates, that $\gamma=\T^1\times\{0\}$. We fix $F$ a lift of $f$ and assume, with no loss in generality, that $\rho(x_0,F)=0$ and that $\rho(x_1,F)=L\ge 1$.

Let $U\subset\A$ be the connected component of $\A\setminus \gamma$ containing
both $x_0$ and $x_1$. As we assumed that $\gamma$ meets both
$U_0,U_1$ it is possible to consider two arcs $j_0,j_1$
contained in $\textrm{cl}(U)=U \cup \gamma$ and an arc $i$ so that:
\begin{enumerate}
\item $j_0\subset U_0$, $j_1\subset U_1$.
\item $j_0$ starts in $x_0$, finishes in $p_0\in\gamma$ and
$j_0\setminus \{p_0\}\subset U$.
\item $j_1$ starts in $x_1$, finishes in $p_1\in\gamma$ and
$j_1\setminus \{p_1\}\subset U$.
\item $i$ is an inessential arc in $\gamma$ joining $p_0,p_1,$ such that, if $\widetilde i$ is a lift of $i$ with an endpoint in $\widetilde p_0$, then $\widetilde i$ is a subarc of the line segment joining $\widetilde p_0$ and $\widetilde p_0+(1,0)$
\end{enumerate}

Set $G=F^2$. Consider the arc $\theta$ given by the concatenation of $j_0\cdot i\cdot (j_1)^{-1}$ and take a lift
$\widetilde{\theta}$ of $\theta$ starting at $\widetilde{x}_0$, a lift of $x_0$, and finishing
at $\widetilde{x}_1$, a lift of $x_1$. As $G(\widetilde{x}_0)=\widetilde{x}_0$ and
$G(\widetilde{x}_1)=\widetilde{x}_1+(2L,0)$ for some integer $L>0$, it holds that $G(\widetilde{\theta})$ meets
$\widetilde{\theta},\widetilde{\theta}+(1,0)$ and $\widetilde{\theta}+(2,0)$. Consider
$\widetilde{j}_0,\widetilde{j}_1,\widetilde{i}$ respective lifts of $j_0, j_1, i$ contained in $\widetilde{\theta}$.
If $\widetilde{p}_0=(a,0)$, and $\widetilde{p}_1=(b,0)$ are the endpoints of $\widetilde{i}$, then $a<b<a+1$. Given two points $(c_1,0), (c_2,0)$, denote by $[c_1, c_2]$ the line segment joining these two points. Let $\psi:\R\to\R$ be such that $G((a,0))=(\psi(a),0)$. Then the following hold

\begin{enumerate}
\item $\psi(a)<a+1$. Indeed, assume for a contradiction this is false. Then, as $j_0\subset U_0$ and $U_0 \cup f^2(U_0)$ is inessential, this implies $G(\tilde{j}_0) \cap
\widetilde{j}_0+(1,0)=\emptyset$, so the point $\widetilde{x}_0+(1,0)$ must be contained in some bounded connected component of $\R^2\setminus \left([a,\psi(a)]\cup\widetilde{j}_0\cup G(\widetilde{j}_0)\right)$. But  $G(\widetilde{j}_0+(1,0))$ connects $\widetilde{x}_0+(1,0)$ to $G((a,0))+(1,0)$ which is contained in the unbounded component of the complement of
    $[a,\psi(a)]\cup\widetilde{j}_0\cup G(\widetilde{j}_0)$, implying that $G(\widetilde{j}_0+(1,0))$
    meets $G(\widetilde{j}_0)\cup\widetilde{j}_0$, which contradicts the definition of a disjoint pair of neighborhoods.

\item $\psi(b)>b-1$: this follows from a similar argument as the one presented for the previous item.
\end{enumerate}

Moreover, it holds:

\begin{enumerate}[resume]

\item $\psi(a)<b$. Indeed, assume for a contradiction this is false. Then $a<b<\psi(a)$, and as $\widetilde{j}_1$
cannot meet $\widetilde{j}_0$ nor $G(\widetilde{j}_0)$, the point $\widetilde{x}_1$ must be contained in some
bounded connected component of $\R^2\setminus \left(\widetilde{j}_0\cup G(\widetilde{j}_0)\cup [a,\psi(a)]\right)$. Hence
$G(\widetilde{j}_1)-(2,0)\cap \R\times\{0\}$ is given by a point contained in $[a,\psi(a)]$ as otherwise
$G(\widetilde{j}_1)-(2,0)$ must intersect $\widetilde{j}_0 \cup G(\widetilde{j}_0)$. Thus
$G(\widetilde{j}_1) \cap \R\times\{0\}$ is given by a point contained in $[a+2,\psi(a)+2]$. Then
$G(\widetilde{i})$ contains both points $G((a,0))$ (due to the first consideration above)
and $(a,0)+(2,0)$ which is impossible as $\widetilde{i}$ lifts an inessential arc in $\gamma$.

\item $\psi(b)>a+2$: the same argument introduced for the case above works
in this one.

\end{enumerate}

But these last two considerations imply that $G(\widetilde{i})$ meets $\widetilde{i}$ and $\widetilde{i}+(2,0)$ which is absurd\textcolor{black}{: $G(\tilde{i})$ would then contain the segment $[a+1,a+2]$, so $f^2(i)$ would be an essential subset of $\A$, contradicting the fact that $f^2$ is a homeomorphism and $i$ is inessential}.
\end{proof}

\end{document}
